\chapter{Predicción con cadenas de Markov visibles}
\label{ch:o2}

Este capítulo construye el primer modelo predictivo del trabajo: una cadena de
Markov de primer orden que, a partir de la celda en la que un individuo se
encuentra hoy, estima en cuál estará mañana. Se llama \emph{visible} porque su
estado (la celda ocupada) es una variable observada directamente, a diferencia de
los estados ocultos que introducirá el capítulo siguiente. Más allá de su interés
propio, el modelo cumple un papel transversal: es la referencia interpretable y
cuantificada que los modelos posteriores (capítulos \ref{ch:o3} y \ref{ch:o4})
deberán superar.

\section{Introducción y objetivo}
\label{sec:o2-intro}

La entrada de este objetivo es la tabla diaria producida en el capítulo
\ref{ch:o1}: una posición por individuo y día natural (UTC), con los huecos
representados de forma explícita y la posición anclada a las 08:00, próxima al
lugar de descanso nocturno (\textit{roost}). Sobre esa secuencia regular se
plantea la pregunta de predicción más simple posible: dado dónde está hoy el ave,
¿dónde estará mañana? El modelo que la responde con menos supuestos, respetando la
estructura temporal de la serie, es una cadena de Markov de primer orden sobre el
espacio discretizado.

La discretización convierte cada posición diaria en un estado: el dominio
geográfico se divide en una rejilla de celdas y la posición del día $t$ se
identifica con la celda que la contiene, $X_t$, dentro de un espacio finito de $N$
celdas. La hipótesis de Markov de primer orden establece que la celda del día
siguiente depende solo de la celda actual, no de toda la trayectoria previa:

\begin{equation}
  P(X_{t+1} = j \mid X_t = i,\, X_{t-1}, \dots, X_0)
  = P(X_{t+1} = j \mid X_t = i) = p_{ij},
  \label{eq:o2-markov}
\end{equation}

donde $p_{ij}$ es la probabilidad de transitar de la celda $i$ a la $j$. El
conjunto de estas probabilidades forma la matriz de transición $P$, con
$\sum_j p_{ij} = 1$ para toda celda de origen. Para captar la estacionalidad del
ciclo migratorio, el modelo no usa una única matriz sino una por mes natural,
$P^{(m)}$ con $m = 1, \dots, 12$. La teoría general de las cadenas de Markov
finitas se trata en el capítulo \ref{ch:preliminares}; aquí se emplea en su forma
aplicada al problema.

La cuestión concreta del capítulo es si una cadena de Markov de primer orden,
global (una sola familia de matrices compartida por todos los individuos) y con
estacionalidad mensual, mejora a la predicción trivial de \emph{persistencia}
(suponer que mañana el ave seguirá en la misma celda que hoy). Esta baseline no es
un rival ingenuo: como se verá, la marcada tendencia de la especie a permanecer en
el mismo lugar la hace sorprendentemente difícil de batir, y esa dificultad es,
por sí misma, un resultado. El modelo de Markov aporta además el criterio de
comparación cuantificado que reutilizan los objetivos siguientes: sus predicciones
se almacenan para el análisis comparativo del error del capítulo \ref{ch:o5}.

\section{Discretización del espacio}
\label{sec:o2-discretizacion}

El primer paso convierte la posición continua en un estado discreto. El dominio
geográfico se cubre con una rejilla de celdas cuadradas de lado \texttt{cell\_deg}
grados, y cada posición $(\texttt{lon}, \texttt{lat})$ se asigna a la celda que la
contiene, identificada por su índice de latitud
$\lfloor \texttt{lat} / \texttt{cell\_deg} \rfloor$ y su índice de longitud
$\lfloor \texttt{lon} / \texttt{cell\_deg} \rfloor$ (por ejemplo, \texttt{122\_49}). Solo se conservan como estados las celdas con al
menos una observación válida, que forman el espacio de celdas activas. Con ello la
posición diaria de un individuo queda reducida a la secuencia de celdas $X_t$ sobre
la que opera la cadena.

El tamaño de celda no es un parámetro inocuo: fija a la vez la resolución espacial
del modelo y la proporción de \emph{self-loops} (transiciones que permanecen en la
misma celda de un día al siguiente), que es justamente la magnitud de la que
depende la fortaleza de la baseline de persistencia. Una celda demasiado grande
hace que casi todos los días parezcan estacionarios y una demasiado pequeña deja
sin datos a la mayoría de las transiciones, así que la elección se decide
comparando cuatro candidatos (0,25\textdegree, 0,5\textdegree, 1\textdegree\ y
2\textdegree; figura \ref{fig:o2-grid}).

El compromiso es claro en los extremos. Con celdas de 2\textdegree\ (del orden de
222 km) el 86,1\,\% de las transiciones son self-loops: el movimiento migratorio
que el modelo debería capturar se diluye dentro de la celda y la persistencia se
infla por construcción. Con celdas de 0,25\textdegree\ (unos 28 km) la resolución
es fina pero la estimación se vuelve inviable, porque en un mes típico cada celda
registra menos de una transición (la mediana de transiciones por celda y mes es
0,86; panel B de la figura), de modo que la mayoría de las celdas se queda sin datos
con los que estimar sus probabilidades. El tamaño intermedio de \textbf{0,5\textdegree}
(del orden de 56 km) es comparable a la escala del grueso de los desplazamientos
diarios (panel A de la figura): un día de movimiento real cruza varias celdas
mientras que un día de descanso permanece en la misma. A esa resolución el espacio
tiene 1\,217 celdas activas y el 73,0\,\% de las transiciones son self-loops, una
residencialidad alta pero no exagerada y coherente con la fenología de la especie.
Se adopta, por tanto, \texttt{cell\_deg} = 0,5\textdegree.

\begin{figure}[H]
  \centering
  \includegraphics[width=\linewidth]{o2_fig01_grid-size-tradeoff.png}
  \caption[Compromiso del tamaño de celda]{Comparación de cuatro tamaños de celda
  para la discretización del espacio (0,25\textdegree, 0,5\textdegree, 1\textdegree\
  y 2\textdegree). Panel A: distribución del desplazamiento diario observado, con
  líneas verticales en el tamaño físico aproximado de cada celda. Panel B: métricas
  de cobertura y densidad por candidato (celdas activas, transiciones, mediana de
  transiciones por celda-mes, porcentaje de pares observados y de self-loops). Panel
  C: cada rejilla superpuesta a la nube de \textit{fixes} válidos. Se adopta
  0,5\textdegree\ porque equilibra resolución espacial y densidad estadística sin
  inflar artificialmente los self-loops.}
  \label{fig:o2-grid}
\end{figure}
